\documentclass[12pt]{article}
\usepackage{amsmath,amssymb,amsthm}
\usepackage{geometry}
\usepackage{hyperref}
\usepackage{natbib}
\geometry{margin=1in}

\newtheorem{definition}{Definition}
\newtheorem{theorem}{Theorem}
\newtheorem{proposition}{Proposition}
\newtheorem{remark}{Remark}
\newtheorem{conjecture}{Conjecture}
\newtheorem{prediction}{Prediction}
\newtheorem{example}{Example}

\title{Behavioral Equivalence Without Representational Equivalence:\\
Japanese Tit (\textit{Parus minor}) Communication\\
as a Case Study in Metric Grounding}

\author{Minamo Minamoto\\
\small Independent Researcher\\
\small ORCID: 0009-0002-1201-5704\\
\small \href{https://github.com/minamominamoto/topology-of-grounding}{topology-of-grounding}
}

\date{April 2026}

\begin{document}
\maketitle

\begin{abstract}
We interpret Suzuki and colleagues' experimental findings on Japanese tit
(\textit{Parus minor}) communication within the framework of metric-geometric
symbol grounding developed in Minamoto (2026).
Two aspects of Japanese tit communication serve as the primary test cases:
(i) the referential grounding of alarm calls to external predator categories
instantiates the formal notion of a grounding map
$G: \Sigma \to \mathcal{F}(S(A))$
assigning call types to distinct regions of sensory space;
(ii) the combinatorial syntax of alarm call sequences
(``ABC'' scan-for-danger notes concatenated with ``D'' approach-caller note
produce a mobbing-like approach response qualitatively distinct
from the response to either element alone)
provides behavioral evidence for structured representational combination.
We derive a quantitative prediction:
RSA \citep{kriegeskorte2008} applied to neural data
from Japanese tit auditory circuits should yield
$d_\mathrm{cog}(\text{Japanese tit}, \text{human}) > 0$---
measured as Gromov--Hausdorff distance between
the metric spaces induced on $\Sigma$ by each agent's grounding map---
even for calls whose referents are behaviourally equivalent across species.
We further discuss the recently discovered wing-fluttering gesture
and its implications for multimodal extensions of the framework.
\end{abstract}

\section{Introduction}

The symbol grounding problem, as originally posed by \citet{harnad1990},
asks whether the symbols manipulated by a cognitive system
are causally connected to sensorimotor states,
or merely formal tokens with no intrinsic meaning.
In Minamoto (2026) we argued that this binary formulation is insufficient:
grounding is not a property a system either has or lacks,
but a \emph{structural} property characterised by the metric geometry
of the sensory space in which grounding occurs.
Two agents that are both weakly grounded may nevertheless
possess representations whose metric geometries are non-isometric,
rendering complete conceptual sharing impossible.

We follow \citet{seyfarth2003} in distinguishing
\emph{functionally referential} calls---those that elicit
context-specific responses in receivers---from
fully \emph{semantically referential} signals in the human linguistic sense.
The Japanese tit's alarm calls are functionally referential;
the question of their deeper semantic status
is precisely the kind of structural question the framework addresses.
We use \emph{behavioral equivalence} to mean:
two agents produce the same pattern of observable responses
to a given call type (e.g.\ both show snake-inspection behaviour
in response to the snake-specific call).

The question of whether combinatorial structure in communication
is specifically human is central to language evolution research
\citep{hockett1960,hauser2002}.
Recent experimental work by Suzuki and colleagues
on Japanese tit (also referred to as Japanese great tit in the primary
literature when classified as a subspecies of \textit{Parus major};
here we follow the independent-species classification \textit{Parus minor})
provides an unusually tractable empirical handle on these questions.
Three findings are relevant:

\begin{enumerate}
\item Japanese tits produce referential alarm calls
      that elicit different predator-searching behaviours
      depending on call type \citep{suzuki2012}.
      Snake-specific calls elicit ground-directed scanning behaviour;
      other call types elicit horizontally-oriented scanning
      appropriate to aerial and mammalian predators \citep{suzuki2012},
      and coordinate multi-species mobbing responses \citep{suzuki2016mob}.

\item Within the general alarm call system,
      ``ABC'' notes (scan-for-danger) concatenated with a ``D'' note
      (approach-the-caller) produce a mobbing-like approach response
      qualitatively distinct from the response to either element alone
      (note: the ``D'' note alone does elicit partial approach behaviour;
      the combined ABC+D response integrates scanning and approach)
      \citep{suzuki2016syn}.
      The reverse ordering ``D\,+\,ABC'' does not elicit
      the same combined response; this was confirmed using
      novel heterospecific sequences \citep{suzuki2017ord}.

\item Japanese tits use wing-fluttering gestures
      to coordinate nest-box entry with their mates \citep{suzuki2024},
      establishing gestural communication as a distinct signalling modality
      in the species.
\end{enumerate}

Together these findings make the Japanese tit
among the most richly characterised non-human vertebrates
for the empirical study of grounding structure.
In this paper we:
(a) show how findings 1 and 2 map onto formal constructs
in the grounding framework,
(b) derive a novel quantitative prediction (Prediction~\ref{pred:rsa}),
and (c) discuss finding 3 and directions for multimodal extensions.

\section{Background: Metric Grounding Framework}

We briefly recall the relevant definitions from Minamoto (2026).
Throughout, each channel activation space $A_i$
is a compact metrizable space
equipped with a specific metric $\rho_i$.

\begin{definition}[Sensory Space]
\label{def:sensory}
For agent $A$ with transduction channels $\{T_1, \ldots, T_n\}$,
the \emph{sensory space} is
$(S(A), \rho_S)$ where $S(A) = \prod_{i=1}^n A_i$
and $\rho_S$ is the $\ell^2$ product metric.
The $\ell^2$ choice is conventional;
conclusions about $d_\mathrm{cog} > 0$ may not be preserved
under independent rescaling of metrics across agents
(see Definition~\ref{def:cogdist}).
\end{definition}

\begin{definition}[Grounding Map and Symbol Distance]
\label{def:groundmap}
Let $\mathcal{F}(S(A))$ denote the hyperspace of
non-empty closed subsets of $S(A)$,
equipped with the Hausdorff metric $d_H$
(which generates the Vietoris topology
when $S(A)$ is compact metric \citep{beer1993}).
A \emph{grounding map} is an injective function
$G: \Sigma \to \mathcal{F}(S(A))$.
When $\Sigma$ is finite (as in the Japanese tit application),
any function from $\Sigma$ is automatically continuous
(since finite sets carry the discrete topology);
the injectivity condition ensures $d_A$ is a genuine metric.
The \emph{symbol distance} induced by $G$ is:
$d_A(\sigma, \sigma') = d_H(G_A(\sigma), G_A(\sigma'))$.
\end{definition}

\begin{remark}[Topological content for finite $\Sigma$]
\label{rem:discrete}
When $\Sigma$ is finite and $G$ is injective,
the metric topology on $(\Sigma, d_A)$ is discrete,
and any two finite sets of the same cardinality are homeomorphic.
The substantive content of the framework
in the finite-$\Sigma$ case lies in the \emph{metric geometry}
of $(\Sigma, d_A)$---the distances between symbols---
rather than in topological invariants.
\end{remark}

\begin{definition}[Cognitive Distance; v2 signature form]
\label{def:cogdist}
Since $|\Sigma|$ is finite, $(\Sigma, d_A)$ and $(\Sigma, d_B)$ are
compact metric spaces. Following the persistence signature
formulation of \cite[Def.~\ref{def:dcog-signature}]{minamoto2026topology},
the \emph{cognitive distance} between agents $A$ and $B$ is the tuple
\[
  d_\mathrm{cog}(A, B)
  = \Big(\,
      d_\mathrm{GH}\bigl((\Sigma, d_A), (\Sigma, d_B)\bigr),\;
      \{d_\mathrm{bn}^{(k)}(A, B)\}_{k \geq 0},\;
      \{\mathrm{PD}_k(A), \mathrm{PD}_k(B)\}_{k \geq 1}\Big),
\]
where $d_\mathrm{bn}^{(k)}$ is the bottleneck distance between
$k$-th Vietoris--Rips persistence diagrams. For the Japanese-tit
prediction developed in this paper, we focus on the $d_\mathrm{GH}$
and $d_\mathrm{bn}^{(0)}$ components, which reflect worst-case
pairwise distortion and merge-structure alignment respectively;
these are the components accessible with the twenty-five--symbol
concept inventory of the tit call system. The $d_\mathrm{bn}^{(k)}$
components for $k \geq 1$ require larger concept inventories to
stabilize (see \cite[App.~\ref{app:dcog-h1}]{minamoto2026topology});
their estimation is left to future work.
\end{definition}

The formal structure of $d_\mathrm{cog}$ decomposes into
a mathematical theorem and a philosophical conjecture:

\begin{theorem}[GH component and isometry; standard result]
\label{thm:gh}
The $d_\mathrm{GH}$ component of $d_\mathrm{cog}(A, B)$
vanishes if and only if $(\Sigma, d_A)$ and $(\Sigma, d_B)$ are
isometric.
\end{theorem}

\begin{remark}[On the signature components]
The full tuple $d_\mathrm{cog}(A, B)$ vanishes (in all components)
if and only if the two metric spaces are isometric; the
$d_\mathrm{bn}^{(k)}$ components vanish in the isometric case
by the stability theorem \cite{cohenSteiner2007}. The isometry
criterion therefore remains the operative equivalence relation for
complete grounding sharing, as in the v1 formulation; v2 adds
intermediate structural invariants that are separately measurable
when complete isometry fails.
\end{remark}

\begin{conjecture}[Metric Non-Isometry Conjecture]
\label{conj:mni}
If $(\Sigma, d_A) \not\cong (\Sigma, d_B)$
(i.e.\ $d_\mathrm{cog}(A,B) > 0$),
then $A$ and $B$ cannot completely share intensional conceptual content.
Proof for the finite-$\Sigma$ case is given in Appendix~\ref{app:proof}.
\end{conjecture}

Theorem~\ref{thm:gh} is a standard result from metric geometry.
Conjecture~\ref{conj:mni} is the substantive non-trivial claim.
The proof (Appendix~\ref{app:proof}) constructs a specific
similarity relation that is preserved under $d_A$ but not under $d_B$
whenever the two spaces are non-isometric.

\section{Japanese Tit Communication as Grounding Structure}

\subsection{Referential Calls: Grounding Map Instantiation}

The referential alarm call system of the Japanese tit \citep{suzuki2012}
establishes a grounding map in which distinct call types
are linked to distinct sensory-behavioural states.
This categorical structure is consistent with the grounding map
assigning elements of $\mathcal{F}(S(P))$ to call types
such that snake and non-snake categories are well-separated
in the symbol distance $d_P$:
\[
d_H(G(\sigma_\mathrm{snake}), G(\sigma_\mathrm{nonsnake})) \;\gg\;
  d_H(G(\sigma_i), G(\sigma_j)) \quad \text{for }
  \sigma_i, \sigma_j \text{ within the same category.}
\]
This is a claim about metric geometry, not topological homeomorphism type
(Remark~\ref{rem:discrete}).

\subsection{Call Concatenation: Structured Representational Combination}

The combinatorial finding of \citet{suzuki2016syn}---
that ``ABC'' notes concatenated with ``D'' note generate
a mobbing-like approach response
not predicted by either element alone---
provides behavioral evidence for representational combination
sensitive to element ordering.
The order-dependence was confirmed using novel heterospecific sequences
\citep{suzuki2017ord}.

In the framework, this is modelled as a Merge operation
$\mathrm{Merge}: \Sigma \times \Sigma \to \Sigma'$
where the metric structure of $\Sigma'$ is not recovered
from that of $\Sigma \times \Sigma$ under the product metric alone:
the combined symbol ``ABC\,+\,D'' has a grounding region
in $\mathcal{F}(S(P))$ that is not determined
by $G(\sigma_\mathrm{ABC})$ and $G(\sigma_\mathrm{D})$ independently.

We note explicitly what remains to be established:
(i) the order-sensitive behavior is consistent with
many computational accounts (associative learning,
Bayesian sequence models, rule-based decoding)
that do not require a grounding-map interpretation;
(ii) for finite $\Sigma$, the notion of ``strictly finer topology''
is not meaningful (all finite metric spaces have discrete topology);
the correct claim is that Merge generates a symbol distance matrix
on $\Sigma'$ not recoverable from the product distance on $\Sigma \times \Sigma$.
We demonstrate this in the following toy example.

\begin{example}[Merge non-recoverability]
\label{ex:merge}
Let $\Sigma = \{\sigma_\mathrm{ABC},\, \sigma_\mathrm{D}\}$
and $\Sigma' = \{\sigma_\mathrm{ABC},\, \sigma_\mathrm{D},\, \sigma_\mathrm{ABC+D}\}$.

Assign symbol distances consistent with the behavioural data:
\begin{align*}
d_P(\sigma_\mathrm{ABC},\, \sigma_\mathrm{D}) &= 1
\end{align*}
(the two calls are distinct and elicit distinct partial responses).
Under the product metric on $\Sigma \times \Sigma$,
the only available distances are derived from $d_P(\sigma_\mathrm{ABC}, \sigma_\mathrm{D}) = 1$:
any reasonable product-metric extension to $\Sigma'$ would predict
that $\sigma_\mathrm{ABC+D}$ lies at distance $\leq 1$ from both components,
and that its distances to each component are determined by
$d_P(\sigma_\mathrm{ABC}, \sigma_\mathrm{D})$ alone.

The behavioural data, however, show that $\sigma_\mathrm{ABC+D}$ elicits
a combined approach-plus-scanning response qualitatively distinct from
either $\sigma_\mathrm{ABC}$ (scan only) or $\sigma_\mathrm{D}$ (approach only).
This means $G_P(\sigma_\mathrm{ABC+D})$ is a new region of sensory-motor
state space not equal to $G_P(\sigma_\mathrm{ABC}) \cup G_P(\sigma_\mathrm{D})$
nor any convex combination thereof.
Consequently:
\[
d_P(\sigma_\mathrm{ABC+D},\, \sigma_\mathrm{ABC})
\;\neq\; d_P(\sigma_\mathrm{ABC+D},\, \sigma_\mathrm{D}),
\]
and neither value is recoverable from $d_P(\sigma_\mathrm{ABC}, \sigma_\mathrm{D})$
under the product metric.
The distance matrix of $(\Sigma', d_P)$ has strictly more structure
than the product distance on $\Sigma \times \Sigma$ can supply.
\end{example}

The grounding-map account offers unified treatment of
(i) referential grounding, (ii) order-sensitive combination, and
(iii) cross-species representational divergence simultaneously.

\section{The Main Prediction}

\begin{prediction}[RSA Discrimination]
\label{pred:rsa}
Let $H$ (human) and $P$ (Japanese tit, \textit{Parus minor})
be agents that have both learned the functional rule
``snake call predicts snake-inspection behaviour.''
Assume a shared symbol set $\Sigma$ (an idealization; see Remark below).
Assuming Conjecture~\ref{conj:mni}:
\[
d_\mathrm{cog}(H, P) > 0.
\]
Operationally: RSA \citep{kriegeskorte2008}
applied to neural responses in auditory cortex ($H$)
and auditory forebrain NCM/CM \citep{mello1992} in ($P$)
to the same set of referential calls
will yield representational geometry matrices
$\mathbf{R}_H, \mathbf{R}_P$ with
$\|\mathbf{R}_H - \mathbf{R}_P\|_F > \delta_0$,
where $\delta_0 > 0$ is pre-specified
and calibrated against within-species baseline (item~4 below).
\end{prediction}

\begin{prediction}[Directional Distance Asymmetry]
\label{pred:asymmetry}
The non-isometry between $(\Sigma, d_P)$ and $(\Sigma, d_H)$
has a specific directional structure, not merely a magnitude.

In the tit ($P$), the snake-specific alarm call $\sigma_\mathrm{snake}$
is associated with a qualitatively distinct sensorimotor response
(ground-directed scanning and vertical retreat)
that has no analogue in the aerial or mammalian predator responses.
This predicts:
\[
d_P(\sigma_\mathrm{snake},\, \sigma_\mathrm{aerial})
\;>\; d_P(\sigma_\mathrm{aerial},\, \sigma_\mathrm{mammal}).
\]
That is, in the tit's representational geometry, the snake call is
more distant from other alarm calls than those alarm calls
are from each other.

In humans ($H$), who have no species-specific predator associations
with the tit alarm calls, the three call types are processed
as acoustically distinct but semantically equivalent unfamiliar signals.
This predicts approximate symmetry:
\[
d_H(\sigma_\mathrm{snake},\, \sigma_\mathrm{aerial})
\;\approx\; d_H(\sigma_\mathrm{aerial},\, \sigma_\mathrm{mammal})
\;\approx\; d_H(\sigma_\mathrm{snake},\, \sigma_\mathrm{mammal}).
\]

The critical test is therefore not only whether
$\|\mathbf{R}_H - \mathbf{R}_P\|_F > \delta_0$
(Prediction~\ref{pred:rsa}),
but whether the \emph{asymmetry}
$d_P(\sigma_\mathrm{snake}, \sigma_\mathrm{aerial}) > d_P(\sigma_\mathrm{aerial}, \sigma_\mathrm{mammal})$
is present in $\mathbf{R}_P$ and absent in $\mathbf{R}_H$.
This directional prediction is specific to the
sensorimotor-grounding account and cannot be derived
from acoustic similarity alone.
In particular, distances tied to survival-relevant categorical
boundaries (predator vs.\ non-predator) are predicted to be
\emph{conserved} as large distances in both species' RSA matrices,
while distances between sequence-structured signals
(e.g.\ $\sigma_\mathrm{ABC}$ vs.\ $\sigma_\mathrm{ABC+D}$)
are expected to \emph{diverge} across species,
since the combinatorial grounding is tit-specific.
\end{prediction}

\begin{remark}[Theoretical justification and the d\_cog/RSA gap]
\label{rem:pred}
The theoretical argument for $d_\mathrm{cog}(H,P) > 0$ is as follows.
Japanese tit alarm calls are processed as conspecific warning signals
coupled to specific sensorimotor contingencies
(ground-scanning for snakes, fleeing from aerial predators, etc.);
humans processing the same acoustic stimuli have no such
species-specific predator associations.
Consequently, the grounding maps $G_P$ and $G_H$
embed the same $\Sigma$ into structurally different regions
of $S(P)$ and $S(H)$ respectively,
yielding $d_P \not\cong d_H$ and hence $d_\mathrm{cog}(H,P) > 0$
under Conjecture~\ref{conj:mni}.
Note that this argument concerns the \emph{auditory-sensorimotor}
grounding of the calls, not visual system differences.

The apparent logical gap between $d_\mathrm{cog}$ and
$\|\mathbf{R}_H - \mathbf{R}_P\|_F$ (Frobenius norm of RSA matrix
differences) noted in v1 of this paper is partially resolved by the
signature formulation of Definition~\ref{def:cogdist}. Under the
signature, the two most immediate components of
$d_\mathrm{cog}(H, P)$---namely $d_\mathrm{GH}$ and
$d_\mathrm{bn}^{(0)}$---are computable directly from the RSA-derived
metric on $\Sigma$. The Frobenius difference
$\|\mathbf{R}_H - \mathbf{R}_P\|_F$ bounds (and is bounded by) the
$d_\mathrm{GH}$ component up to constants that depend on $|\Sigma|$;
the $d_\mathrm{bn}^{(0)}$ component is stable under RSA matrix
perturbation by the Cohen--Steiner stability theorem
\cite{cohenSteiner2007}. The RSA operationalization therefore tracks
well-defined signature components, not a single scalar
$d_\mathrm{cog}$; Prediction~\ref{pred:rsa} below should be read as
predicting non-vanishing of these tractable components, with the
$d_\mathrm{bn}^{(k)}$ components for $k \geq 1$ requiring larger
concept inventories than the twenty-five--symbol tit system
provides.
\end{remark}

\begin{remark}[Shared-$\Sigma$ assumption]
The prediction assumes a shared symbol set $\Sigma$,
but the grounding of each call type is species-specific.
Humans process the tit calls as unfamiliar acoustic stimuli,
while tits process them as conspecific alarm signals.
This non-equivalence means the measured $\mathbf{R}_H$
reflects acoustic processing without predator-linked associations,
while $\mathbf{R}_P$ reflects semantically loaded processing.
The prediction is that even under this asymmetry,
$d_\mathrm{cog}(H,P) > 0$ is detectable via RSA.
\end{remark}

A positive result supports Prediction~\ref{pred:rsa};
a negative result would not definitively refute it
given methodological challenges (modality mismatch between
electrophysiology and fMRI/MEG, SNR differences, temporal resolution).
Technical limitations are discussed in Minamoto (2026, Section~5).

\subsection{Experimental Design}

\begin{enumerate}
\item Neural recordings from Japanese tit auditory forebrain (NCM/CM)
      during playback of the full call repertoire.
      Note: Mello et al. \citeyearpar{mello1992} identified NCM
      in the zebra finch (\textit{Taeniopygia guttata}); the region is
      considered homologous across oscine songbirds
      (Avian Brain Nomenclature Consortium; \citealt{jarvis2005}).
\item fMRI/MEG recordings from human participants
      hearing the same call stimuli.
\item RSA \citep{kriegeskorte2008} between the two neural geometry matrices.
\item Comparison with RSA baseline from two Japanese tit individuals.
      The theory predicts $d_\mathrm{cog}(P_1, P_2) \ll d_\mathrm{cog}(H, P)$
      for individuals with comparable ecological histories,
      call repertoire experience, and predator exposure---
      that is, within-species cognitive distance should be
      significantly smaller than cross-species distance,
      not necessarily zero.
      $\delta_0$ is set as a fixed multiple of within-species variance,
      pre-registered prior to data collection.
\end{enumerate}

\section{Proof of Concept: Cross-Agent RSA Discrimination}
\label{sec:pilot}

The experimental design of Section~4 requires cross-species RSA
between Japanese tit auditory forebrain and human auditory cortex.
Before that experiment can be conducted, it is essential to verify
that the RSA-based $\rho_\mathrm{cog}$ protocol actually
discriminates structurally distinct representational geometries.
We report a proof-of-concept using language model checkpoints,
which constitute agents with known structural differences.

\paragraph{Design rationale.}
The cross-species comparison (human vs.\ \textit{Parus minor})
is structurally analogous to a cross-architecture comparison:
both compare RSA dissimilarity matrices from agents that process
the same stimuli through different grounding structures.
If $\rho_\mathrm{cog}$ fails to discriminate cross-architecture
differences in language models---where ground truth is available---
it would provide little confidence for the cross-species case.

\paragraph{Procedure.}
We computed RSA dissimilarity matrices over $n = 29$ probe concepts
at hidden layers $\ell \in \{2, 4, 6\}$ of two architecturally
distinct models: Pythia-70m \citep{biderman2023} (fully trained,
step~143000) and GPT-2 small \citep{radford2019language}.
Three prompt templates were averaged to reduce surface-form bias.
As a control, we compared GPT-2 with Pythia at random
initialization (step~0), where architectural differences are
unconfounded by training.

\paragraph{Results.}
Across all layers:
\begin{itemize}
  \item $\rho_\mathrm{cog}(\text{GPT-2}, \text{Pythia-trained})
    \approx 0.42$ (trained agents sharing training-induced geometry)
  \item $\rho_\mathrm{cog}(\text{GPT-2}, \text{Pythia-step-0})
    \approx 0.10$ (trained vs.\ random: structurally distinct)
\end{itemize}
The protocol reliably discriminates between structurally similar
and structurally distinct agents, with a fourfold ratio between
the two conditions.

\paragraph{Implications for the tit experiment.}
This result establishes that $\rho_\mathrm{cog}$ has sufficient
sensitivity to detect structural differences between agents
sharing the same stimulus set but processed through different
representational geometries---precisely the situation in the
human vs.\ \textit{Parus minor} comparison.
The expected cross-species $\rho_\mathrm{cog}(H, P)$ under
the sensorimotor-grounding account should be substantially
lower than within-species baseline $\rho_\mathrm{cog}(P_1, P_2)$,
analogous to the trained vs.\ random contrast observed here.
The within-species baseline in step~4 of the experimental design
corresponds to the Pythia-trained vs.\ Pythia-trained control.
All code is available at
\url{https://github.com/minamominamoto/grounding-without-reality}.

\section{Discussion}

\subsection{Wing-Fluttering Gestures and Multimodal Extensions}

The discovery of wing-fluttering gestures in Japanese tits \citep{suzuki2024}
suggests extending the framework to settings where
auditory and visual channels interact in a context-dependent manner.
The ``after you'' coordination function of this gesture
motivates, in principle, a fiber bundle model of the sensory space.
Any fiber bundle over a finite discrete context space
is trivially trivializable; non-trivial bundle structure
would require the context space to carry continuous topological structure,
which remains to be formalised from the present data.
This is an open problem (Open Problem 2 of Minamoto 2026).

\subsection{Combinatorial Structure in Animal Communication}

The question of whether combinatorial structure in signalling
is specifically human is central to language evolution research
\citep{hockett1960,hauser2002}.
The Japanese tit data are consistent with the hypothesis
that combinatorial grounding structure---
meaning generated by ordered combination
rather than by individual elements---
represents a convergent solution under predator-diversity pressure.
This hypothesis requires comparative data from independent lineages
to be tested; the present paper offers one consistent case.

\section{Conclusion}

Japanese tit communication research provides
an empirical anchor for the metric grounding framework.
The referential alarm call system \citep{suzuki2012,suzuki2016mob}
is consistent with the formal grounding map $G: \Sigma \to \mathcal{F}(S(A))$
assigning call types to distinct regions of sensory space.
Call concatenation \citep{suzuki2016syn,suzuki2017ord}
provides behavioral evidence for metric combination structure
not recoverable from the product of individual call geometries.
The main quantitative prediction---
$d_\mathrm{cog}(\text{human}, \text{Japanese tit}) > 0$
conditional on Conjecture~\ref{conj:mni}---
is directly testable via RSA \citep{kriegeskorte2008},
subject to the d$_\mathrm{cog}$/RSA approximation noted in
Remark~\ref{rem:pred}.

\appendix
\section{Behavioral Equivalence Does Not Imply Metric Equivalence}
\label{app:proposition}

\begin{proposition}[Behavioral Equivalence Does Not Imply Metric Equivalence]
\label{prop:main}
Let $\Sigma$ be a finite set of symbols,
and let two agents $A$ and $B$ share an identical
stimulus-response mapping $f_A = f_B : \Sigma \to \mathcal{R}$.
There exist grounding maps $G_A, G_B$ inducing metric spaces
$(\Sigma, d_A)$ and $(\Sigma, d_B)$ such that:
\begin{enumerate}
  \item $f_A = f_B$ \emph{(behavioral equivalence)},
  \item $(\Sigma, d_A)$ and $(\Sigma, d_B)$ are \emph{not isometric}.
\end{enumerate}
\end{proposition}

\begin{proof}
Let $\Sigma = \{s_1, s_2, s_3\}$ with identical response mapping
$f(s_i) = r_i$ for $i = 1, 2, 3$.

Define $d_A$ (collinear geometry):
\[
d_A(s_1, s_2) = 1,\quad
d_A(s_2, s_3) = 1,\quad
d_A(s_1, s_3) = 2.
\]

Define $d_B$ (equilateral geometry):
\[
d_B(s_1, s_2) = 1,\quad
d_B(s_2, s_3) = 1,\quad
d_B(s_1, s_3) = 1.
\]

Both $d_A$ and $d_B$ are valid metrics on $\Sigma$.
Both preserve the same stimulus-response mapping $f$.
Yet $(\Sigma, d_A)$ and $(\Sigma, d_B)$ are not isometric:
in $(\Sigma, d_A)$, the pair $(s_1, s_3)$ achieves the unique
maximum distance 2, while in $(\Sigma, d_B)$ all pairs are
equidistant at 1.
No bijection $\phi: \Sigma \to \Sigma$ can simultaneously preserve
all pairwise distances, so the spaces are non-isometric.
Therefore behavioral equivalence ($f_A = f_B$) does not determine
representational geometry.
\end{proof}

\begin{remark}[Connection to the Japanese tit]
In the Japanese tit application, $s_1, s_2, s_3$ correspond to
$\sigma_\mathrm{snake}, \sigma_\mathrm{aerial}, \sigma_\mathrm{mammal}$.
The tit and the human share the behavioral mapping in the sense that
both can be trained to respond differentially to the three call types.
Proposition~\ref{prop:main} establishes that this shared mapping
does not determine whether the underlying representational geometries
are equivalent---which is precisely the claim of
Conjecture~\ref{conj:mni} and the content of
Predictions~\ref{pred:rsa}--\ref{pred:asymmetry}.
\end{remark}

\section{Proof of Conjecture~\ref{conj:mni} for Finite Symbol Sets}
\label{app:proof}

We prove that non-isometry of finite metric spaces implies the
existence of a similarity relation present in one space but not the other,
constituting a conceptual content that cannot be completely shared.

\paragraph{Setup.}
Let $|\Sigma| = n$ and let $(\Sigma, d_A)$ and $(\Sigma, d_B)$
be finite metric spaces on the same symbol set.
Define the \emph{similarity ordering} of agent $A$ as the preorder
$\preceq_A$ on unordered pairs $\{{\sigma, \sigma'}\} \subseteq \Sigma$:
\[
\{\sigma, \sigma'\} \preceq_A \{\tau, \tau'\}
\iff d_A(\sigma, \sigma') \leq d_A(\tau, \tau').
\]
This ordering encodes all pairwise similarity relations derivable
from the grounding map: $\sigma$ and $\sigma'$ are at least as similar
as $\tau$ and $\tau'$ iff their grounding regions are at least as close.

\paragraph{Definition (complete intensional sharing).}
Agents $A$ and $B$ \emph{completely share intensional conceptual content}
with respect to $\Sigma$ if there exists a bijection $\phi: \Sigma \to \Sigma$
such that for all pairs $\{\sigma, \sigma'\}$ and $\{\tau, \tau'\}$:
\[
\{\sigma, \sigma'\} \preceq_A \{\tau, \tau'\}
\iff \{\phi(\sigma), \phi(\sigma')\} \preceq_B \{\phi(\tau), \phi(\tau')\}.
\]
That is, every similarity ordering relation that $A$ represents
is preserved under $\phi$ in $B$, and vice versa.

\paragraph{Lemma.}
For finite metric spaces, complete intensional sharing (as defined above)
holds if and only if $(\Sigma, d_A)$ and $(\Sigma, d_B)$ are isometric.

\paragraph{Proof.}
($\Rightarrow$) Suppose complete intensional sharing holds via bijection $\phi$.
For any $\sigma, \sigma'$, the strict ordering
$\{\sigma, \sigma'\} \prec_A \{\tau, \tau'\}$ iff
$\{\phi(\sigma), \phi(\sigma')\} \prec_B \{\phi(\tau), \phi(\tau')\}$
holds for all pairs by assumption.
Since $\Sigma$ is finite, the full distance matrix $[d_A(\sigma_i, \sigma_j)]$
is determined by the similarity ordering (up to a common positive scale
if the ordering is total; exactly, if the ordering distinguishes all pairs).
Therefore $\phi$ preserves the metric structure, making the spaces isometric.

($\Leftarrow$) Suppose $(\Sigma, d_A)$ and $(\Sigma, d_B)$ are not isometric.
Then for every bijection $\phi: \Sigma \to \Sigma$,
there exists at least one pair of pairs such that
$d_A(\sigma, \sigma') < d_A(\tau, \tau')$ but
$d_B(\phi(\sigma), \phi(\sigma')) \geq d_B(\phi(\tau), \phi(\tau'))$,
or vice versa.
This pair constitutes a similarity relation represented by $A$
(``$\sigma$ is more similar to $\sigma'$ than $\tau$ is to $\tau'$'')
that $B$ does not represent under $\phi$ in the same direction.
Since $\phi$ was arbitrary, no bijection achieves complete sharing.
Therefore $A$ and $B$ cannot completely share intensional conceptual
content. $\square$

\paragraph{Application to Japanese tit.}
Concretely: the tit's alarm call system gives
$d_P(\sigma_\mathrm{snake}, \sigma_\mathrm{aerial}) >
d_P(\sigma_\mathrm{aerial}, \sigma_\mathrm{mammal})$
(Prediction~\ref{pred:asymmetry}).
If humans process the same calls with
$d_H(\sigma_\mathrm{snake}, \sigma_\mathrm{aerial}) \approx
d_H(\sigma_\mathrm{aerial}, \sigma_\mathrm{mammal})$,
the two metric spaces are not isometric.
By the Lemma, the similarity relation
``snake call is more distinct from aerial calls
than aerial calls are from mammal calls''
is represented in the tit's grounding structure and not in the human's.
This is the conceptual content that cannot be completely shared.

\section*{Acknowledgements}

The author thanks the topology-of-grounding project
(\url{https://github.com/minamominamoto/topology-of-grounding})
for the theoretical framework developed in the companion papers.

\bibliographystyle{plainnat}
\begin{thebibliography}{10}

\bibitem[Beer(1993)]{beer1993}
Beer, G. (1993).
\textit{Topologies on Closed and Closed Convex Sets}.
Dordrecht: Kluwer Academic Publishers.
\href{https://doi.org/10.1007/978-94-015-8149-3}{doi:10.1007/978-94-015-8149-3}

\bibitem[Harnad(1990)]{harnad1990}
Harnad, S. (1990).
The symbol grounding problem.
\textit{Physica D}, 42(1--3), 335--346.
\href{https://doi.org/10.1016/0167-2789(90)90087-6}{doi:10.1016/0167-2789(90)90087-6}

\bibitem[Jarvis et al.(2005)]{jarvis2005}
Jarvis, E.~D., G{\"u}nt{\"u}rk{\"u}n, O., Bruce, L., Csillag, A.,
Karten, H., Kuenzel, W., \ldots\ Butler, A.~B. (2005).
Avian brains and a new understanding of vertebrate brain evolution.
\textit{Nature Reviews Neuroscience}, 6(2), 151--159.
\href{https://doi.org/10.1038/nrn1606}{doi:10.1038/nrn1606}

\bibitem[Hauser et al.(2002)]{hauser2002}
Hauser, M.~D., Chomsky, N., \& Fitch, W.~T. (2002).
The faculty of language: what is it, who has it, and how did it evolve?
\textit{Science}, 298(5598), 1569--1579.
\href{https://doi.org/10.1126/science.298.5598.1569}{doi:10.1126/science.298.5598.1569}

\bibitem[Hockett(1960)]{hockett1960}
Hockett, C.~F. (1960).
The origin of speech.
\textit{Scientific American}, 203(3), 88--96.
\href{https://doi.org/10.1038/scientificamerican0960-88}{doi:10.1038/scientificamerican0960-88}

\bibitem[Kriegeskorte et al.(2008)]{kriegeskorte2008}
Kriegeskorte, N., Mur, M., \& Bandettini, P. (2008).
Representational similarity analysis---connecting the branches
of systems neuroscience.
\textit{Frontiers in Systems Neuroscience}, 2, 4.
\href{https://doi.org/10.3389/neuro.06.004.2008}{doi:10.3389/neuro.06.004.2008}
\bibitem[Cohen-Steiner et al.(2007)]{cohenSteiner2007}
Cohen-Steiner, D., Edelsbrunner, H., \& Harer, J. (2007).
Stability of persistence diagrams.
\textit{Discrete \& Computational Geometry}, 37(1), 103--120.


\bibitem[Mello et al.(1992)]{mello1992}
Mello, C.~V., Vicario, D.~S., \& Clayton, D.~F. (1992).
Song presentation induces gene expression in the songbird forebrain.
\textit{Proceedings of the National Academy of Sciences}, 89(15), 6818--6822.
\href{https://doi.org/10.1073/pnas.89.15.6818}{doi:10.1073/pnas.89.15.6818}

\bibitem[Minamoto(2026)]{minamoto2026topology}
Minamoto, M. (2026a).
On the topology of symbol grounding:
sensory space structure as a determinant of cognitive representation.
Preprint.
\url{https://github.com/minamominamoto/topology-of-grounding}

\bibitem[Seyfarth \& Cheney(2003)]{seyfarth2003}
Seyfarth, R.~M., \& Cheney, D.~L. (2003).
Signalers and receivers in animal communication.
\textit{Annual Review of Psychology}, 54, 145--173.
\href{https://doi.org/10.1146/annurev.psych.54.101601.145121}{doi:10.1146/annurev.psych.54.101601.145121}

\bibitem[Suzuki(2012)]{suzuki2012}
Suzuki, T.~N. (2012).
Referential mobbing calls elicit different predator-searching behaviours
in Japanese great tits.
\textit{Animal Behaviour}, 84(1), 53--57.
\href{https://doi.org/10.1016/j.anbehav.2012.03.030}{doi:10.1016/j.anbehav.2012.03.030}

\bibitem[Suzuki(2016a)]{suzuki2016mob}
Suzuki, T.~N. (2016a).
Referential calls coordinate multi-species mobbing
in a forest bird community.
\textit{Journal of Ethology}, 34, 79--84.
\href{https://doi.org/10.1007/s10164-015-0449-1}{doi:10.1007/s10164-015-0449-1}

\bibitem[Suzuki et al.(2016b)]{suzuki2016syn}
Suzuki, T.~N., Wheatcroft, D., \& Griesser, M. (2016b).
Experimental evidence for compositional syntax in bird calls.
\textit{Nature Communications}, 7, 10986.
\href{https://doi.org/10.1038/ncomms10986}{doi:10.1038/ncomms10986}

\bibitem[Suzuki et al.(2017)]{suzuki2017ord}
Suzuki, T.~N., Wheatcroft, D., \& Griesser, M. (2017).
Wild birds use an ordering rule to decode novel call sequences.
\textit{Current Biology}, 27, 2331--2336.
\href{https://doi.org/10.1016/j.cub.2017.06.031}{doi:10.1016/j.cub.2017.06.031}

\bibitem[Suzuki \& Sugita(2024)]{suzuki2024}
Suzuki, T.~N., \& Sugita, N. (2024).
The `after you' gesture in a bird.
\textit{Current Biology}, 34, R231--R232.
\href{https://doi.org/10.1016/j.cub.2024.01.007}{doi:10.1016/j.cub.2024.01.007}

\bibitem[Biderman et al.(2023)]{biderman2023}
Biderman, S., et al.\ 2023.
Pythia: A suite for analyzing large language models across training and scaling.
\textit{Proceedings of ICML 2023}.

\bibitem[Radford et al.(2019)]{radford2019language}
Radford, A., et al.\ 2019.
Language models are unsupervised multitask learners.
\textit{OpenAI Blog}, 1(8).

\end{thebibliography}

\end{document}
